\section*{To the Student}

In a typical high school math text, each section has a ``technique'' which you practice in a series of exercises very like the examples in the text.  This book is different.  In this course you will be learning to use calculus both as a tool and as a language in which you can think coherently about the problems you will be studying.  As with any other language, a certain amount of time will need to be spent learning and practicing the formal rules.  For instance, the conjugation of {\em \^etre\/} must be almost second nature to you if you are to be able to read a novel---or even a newspaper---in French.  In calculus, too, there are a number of manipulations which must become automatic so that you can focus clearly on the content of what is being said.  It is important to realize, however, that becoming good at these manipulations is not the goal of learning calculus any more than becoming good at declensions and conjugations is the goal of learning French.

Up to now, most of the problems you have met in math classes have had definite answers such as ``17,'' or ``the circle with radius 1.75 and center at (2,3).''  Such definite answers are satisfying (and even comforting).  However, many interesting and important questions, like ``How far is it to the planet Pluto,'' or ``How many people are there with sickle-cell anemia,'' or ``What are the solutions to the equation $x^5 + x + 1 = 0$'' can't be answered exactly.  Instead, we have ways to {\bf approximate} the answers, and the more time and/or money we are willing to expend, the better our approximations may be.  While many calculus problems do have exact answers, such problems often tend to be special or atypical in some way.  Therefore, while you will be learning how to deal with these ``nice'' problems, you will also be developing ways of making good approximations to the solutions of the less well-behaved (and more common!) problems.

The computer or the graphing calculator is a tool that that you will need for this course, along with a clear head and a willing hand.  We don't assume that you know anything about this technology ahead of time.  Everything necessary is covered completely as we go along.

You can't learn mathematics simply by reading or watching others.  The only way you can internalize the material is to work on problems yourself.  It is by grappling with the problems that you will come to see what it is you do understand, and to see where your understanding is incomplete or fuzzy.

One of the most important intellectual skills you can develop is that of exploring questions on your own.  Don't simply shut your mind down when you come to the end of an assigned problem.  These problems have been designed not so much to capture the essence of calculus as to prod your thinking, to get you wondering about the concepts being explored.  See if you can think up and answer variations on the problem.  Does the problem suggest other questions?  The ability to ask good questions of your own is at least as important as being able to answer questions posed by others.

We encourage you to work with others on the exercises.  Two or three of you of roughly equal ability working on a problem will often accomplish much more than would any of you working alone.  You will stimulate one another's imaginations, combine differing insights into a greater whole, and keep up each other's spirits in the frustrating times.  This is particularly effective if you first spend time individually working on the material.  Many students find it helpful to schedule a regular time to get together to work on problems.

 Above all, take time to pause and admire the beauty and power of what you are learning.  Aside from its utility, calculus is one of the most elegant and richly structured creations of the human mind and deserves to be profoundly admired on those grounds alone.  Enjoy!

